\documentclass[11pt,a4paper]{article}
\usepackage[T1]{fontenc}
\usepackage[utf8]{inputenc}
\usepackage{mathptmx}
\usepackage{amsmath,amssymb,amsthm}
\usepackage{graphicx}
\usepackage{booktabs}
\usepackage{xcolor}
\usepackage[margin=2.6cm]{geometry}
\usepackage{caption}
\usepackage{framed}
\usepackage[colorlinks=true,linkcolor=blue!55!black,citecolor=blue!55!black,
            urlcolor=blue!55!black]{hyperref}
\usepackage{fancyhdr}
\usepackage{enumitem}

\captionsetup{font=small,labelfont=bf,width=0.94\textwidth}
\pagestyle{fancy}\fancyhf{}
\fancyhead[L]{\small Power-law suppression and the status of radion leakage}
\fancyhead[R]{\small\thepage}
\renewcommand{\headrulewidth}{0.3pt}

\newcommand{\Mpl}{M_{\rm Pl}}
\newcommand{\hMpc}{h\,\mathrm{Mpc}^{-1}}
\newcommand{\LCDM}{\Lambda\mathrm{CDM}}
\newcommand{\Geff}{G_{\rm eff}}
\newcommand{\meff}{m_{\rm eff}}

\definecolor{boxgrey}{gray}{0.95}
\newenvironment{keybox}
  {\begin{framed}\small\noindent}
  {\end{framed}}

\title{\bfseries A phenomenological power-law suppression\\[2pt]
of small-scale power:\\[6pt]
\large observational status, and why the published\\
radion Lagrangian cannot produce it}

\author{Andre Swart\\[2pt]
\small GeometricCosmo, Cape Town, South Africa\\
\small \texttt{geometriccosmo.illusion559@passinbox.com}}

\date{September 2026 \quad|\quad Version 2.0}

\begin{document}
\maketitle
\thispagestyle{empty}

\begin{abstract}
\noindent
A broken power-law transfer function, $T(k)=(1.5/k)^{1/2}$ for $k\ge1.5\,\hMpc$ and
$T=1$ below, was previously reported to satisfy $S_8$ and Lyman-$\alpha$ constraints
simultaneously while leaving $\sigma_8$ too high, and was attributed to a radion
degree of freedom in a five-dimensional Gauss--Bonnet braneworld. We freeze that
function and its reported numbers as an immutable phenomenological baseline, then
test the baseline and the attribution independently.

The baseline does not survive. An independent \textsc{camb} run at fixed $A_s$ shows
that the function changes $\sigma_8$ by $-4.8\times10^{-5}$ in fractional terms,
because only $0.048\%$ of the $\sigma_8$ variance lies above the break scale.
It follows that $S_8$ is unchanged as well: $S_8$ and $\sigma_8$ differ only by the
fixed factor $(\Omega_m/0.3)^{1/2}=1.025$, so the reported pattern of an $S_8$ pass
alongside a $\sigma_8$ failure is not attainable. Both sit $1.7\sigma$ above the
weak-lensing target. The reported Lyman-$\alpha$ ratio of $0.891$ is reproduced only
if $T$ multiplies $P(k)$ rather than the density contrast, and only if the average
includes two sampling points at which $T\equiv1$ by construction; under the standard
convention the suppression reaches $70\%$ at $k=5\,\hMpc$, which is excluded.

We then derive, from the published radion equation of motion, the quasi-static
effective gravitational coupling
$\Geff/G = 1+2\beta^2k^2/(k^2+a^2\meff^2)$. This fails in three independent ways.
The correction scales as $\beta^2$ and is therefore positive: the mechanism enhances
clustering and cannot lower $\sigma_8$. Its shape is a saturating Yukawa step, not a
declining power law. And because the coupling reaches non-relativistic matter only
through electromagnetic binding energy, the strength required for a percent-level
effect exceeds existing Solar-System and equivalence-principle limits by at least
four orders of magnitude. A previously reported screening factor of $10^{42}$ is
withdrawn as a category error; the correct statement is that the Lagrangian contains
no screening mechanism at all.

We state the conditions any successor mechanism must satisfy, and argue that the
proposed move to six-dimensional flux compactification should be deferred, since
none of the three failures depends on the number of dimensions.
\end{abstract}

\vspace{6pt}
\begin{keybox}
\textbf{Status of this document.} This is a negative result. It supersedes the
observational claims of versions 1.8.x and 1.9.x of this project, which should not
be cited for the conclusions listed in Appendix~D. The 5D background numerics and
the reproducibility infrastructure are unaffected and remain available.
\end{keybox}

\newpage
\tableofcontents
\newpage

%=====================================================================
\section{Introduction}

The amplitude of late-time matter clustering inferred from weak gravitational
lensing sits persistently below the value extrapolated from the cosmic microwave
background under $\LCDM$. Depending on the survey and the treatment of systematics,
weak-lensing analyses prefer $S_8\equiv\sigma_8(\Omega_m/0.3)^{1/2}$ near $0.76$--$0.78$,
while \emph{Planck} 2018 gives $S_8=0.832\pm0.013$. The discrepancy is at the level
of two to three standard deviations and has motivated a large number of proposals
that suppress small-scale power or slow late-time growth.

Any such proposal has to thread a narrow gap. Suppression strong enough to move
$S_8$ must act at scales the $\sigma_8$ filter can see, roughly
$k\lesssim1\,\hMpc$. Suppression that extends to $k\sim1$--$10\,\hMpc$ runs into the
Lyman-$\alpha$ forest, which constrains the one-dimensional flux power spectrum and
therefore the three-dimensional matter power over precisely that range. These two
requirements pull in opposite directions across roughly one decade in wavenumber,
and most mechanisms fail one of them.

Earlier versions of this project proposed that a radion in a five-dimensional
Gauss--Bonnet braneworld, coupled to the electromagnetic sector, produces a broken
power-law transfer function that threads this gap. The claim had two parts, which
we will keep carefully separate:

\begin{description}[leftmargin=2.2em,style=nextline,itemsep=2pt]
\item[Layer 1 --- Phenomenology.] A particular function $T(k)$, applied to the
$\LCDM$ power spectrum, yields a specified set of observables.
\item[Layer 2 --- Mechanism.] Radion leakage in a braneworld produces that function.
\item[Layer 3 --- Derivation.] The higher-dimensional field equations require that
function uniquely.
\end{description}

\noindent
Layer 3 was never claimed. This paper tests Layers 1 and 2. Both fail, Layer 1 for
reasons that have nothing to do with braneworlds and everything to do with the
$\sigma_8$ window function, and Layer 2 for reasons that have nothing to do with the
number of extra dimensions.

Section~\ref{sec:baseline} freezes the baseline. Section~\ref{sec:verify} reports an
independent verification and three corrections. Section~\ref{sec:geff} derives
$\Geff(k,a)$ from the published Lagrangian, and Section~\ref{sec:failures} states the
three ways it fails. Section~\ref{sec:screening} withdraws the earlier screening
result. Sections~\ref{sec:comoving}--\ref{sec:6d} set out what a successor mechanism
would need and why the six-dimensional programme should wait.

%=====================================================================
\section{The frozen baseline (Layer 1)}\label{sec:baseline}

To prevent the target from moving during the analysis, we fix the following and do
not alter them anywhere in this paper.

\paragraph{Transfer function.}
\begin{equation}
T(k)=
\begin{cases}
1, & k<k_c,\\[2pt]
\left(k_c/k\right)^{1/2}, & k\ge k_c,
\end{cases}
\qquad k_c = 1.5\,\hMpc .
\label{eq:Tk}
\end{equation}
Both $k_c$ and the index $1/2$ are fitted quantities. They are not derived from any
five- or six-dimensional calculation, and we do not treat them as such.

\paragraph{Cosmology.} \emph{Planck} 2018 TT,TE,EE$+$lowE$+$lensing best fit, with
$A_s$ held fixed. No amplitude rescaling is permitted at any stage, since allowing
$A_s$ to float would absorb precisely the effect under test. Values are listed in
Table~\ref{tab:cosmo}.

\begin{table}[htb]
\centering
\small
\begin{tabular}{lll}
\toprule
Parameter & Value & Note\\
\midrule
$H_0$        & $67.36$ km s$^{-1}$Mpc$^{-1}$ & \\
$\Omega_b h^2$ & $0.02237$ & \\
$\Omega_c h^2$ & $0.1200$  & \\
$\sum m_\nu$ & $0.06$ eV & one massive eigenstate\\
$\tau$       & $0.0544$  & \\
$A_s$        & $2.100\times10^{-9}$ & \textbf{fixed, never rescaled}\\
$n_s$        & $0.9649$  & \\
$\Omega_k$   & $0$       & \\
\midrule
$\Omega_m$   & $0.31519$ & derived\\
$(\Omega_m/0.3)^{1/2}$ & $1.02501$ & derived; fixes $S_8/\sigma_8$\\
\bottomrule
\end{tabular}
\caption{Frozen cosmological inputs. Linear matter power only; no non-linear
prescription is applied, since the baseline $T(k)$ is a linear-theory object.}
\label{tab:cosmo}
\end{table}

\paragraph{Observables.} With $W(x)=3(\sin x-x\cos x)/x^3$ the real-space top-hat
window,
\begin{equation}
\sigma_R^2=\int\frac{\mathrm{d}k}{k}\,\frac{k^3P(k)}{2\pi^2}\,W^2(kR),
\qquad R=8\,h^{-1}\mathrm{Mpc},
\label{eq:sigma8}
\end{equation}
\begin{equation}
S_8=\sigma_8\left(\Omega_m/0.3\right)^{1/2},
\label{eq:S8}
\end{equation}
and the Lyman-$\alpha$ ratio as originally defined, namely the unweighted mean of
$P_{\rm mod}(k)/P_{\LCDM}(k)$ at $k=0.5,\,1.0,\,2.0,\,3.0\,\hMpc$.

\paragraph{Reported numbers.} The values to be tested are
\begin{equation}
S_8=0.829,\qquad \sigma_8=0.811,\qquad
\text{Lyman-}\alpha\text{ ratio}=0.891,
\end{equation}
previously annotated as \textsc{pass}, \textsc{fail}, \textsc{pass} respectively
against targets $S_8=0.78\pm0.03$, $\sigma_8=0.76\pm0.03$, and a Lyman-$\alpha$
ratio in $[0.70,0.90]$.

%=====================================================================
\section{Independent verification, and three corrections}\label{sec:verify}

We ran \textsc{camb} 2.0.4 with the parameters of Table~\ref{tab:cosmo}, extracted
the linear $P(k)$ at $z=0$ over $10^{-4}\le k\le100\,\hMpc$ on $4000$ logarithmic
points, and evaluated Eqs.~\eqref{eq:sigma8}--\eqref{eq:S8} by Simpson quadrature.
As a check on the quadrature, our $\LCDM$ value $\sigma_8=0.81125$ agrees with
\textsc{camb}'s internal $\sigma_8=0.81107$ to $2\times10^{-4}$ fractionally.
The script is \texttt{verify\_baseline.py}; its output is
\texttt{baseline\_verification.json}.

Three separate problems emerge.

%---------------------------------------------------------------------
\subsection{Correction 1: the transfer function does not move \texorpdfstring{$\sigma_8$}{sigma8} at all}

Applying Eq.~\eqref{eq:Tk} in the standard way, $P_{\rm mod}(k)=T^2(k)P_{\LCDM}(k)$,
gives
\begin{equation}
\sigma_8^{\rm mod}=0.81121,
\qquad
\frac{\sigma_8^{\rm mod}}{\sigma_8^{\LCDM}}-1 = -4.8\times10^{-5}.
\end{equation}
The change is five orders of magnitude smaller than the $6.7\%$ shift that would be
needed to reach $\sigma_8=0.76$.

The reason is visible in Figure~\ref{fig:window}. The integrand of
Eq.~\eqref{eq:sigma8} peaks near $k\simeq0.2\,\hMpc$ and is already negligible by
$k\simeq1\,\hMpc$, because the top-hat window oscillates and decays as $k^{-2}$
beyond $kR\sim\pi$. Integrating the measured integrand, the fraction of the total
$\sigma_8$ variance carried by $k\ge k_c=1.5\,\hMpc$ is
\begin{equation}
\frac{\int_{k_c}^{\infty}\mathrm{d}\ln k\; \mathrm{d}\sigma_8^2/\mathrm{d}\ln k}
     {\int_{0}^{\infty}\mathrm{d}\ln k\; \mathrm{d}\sigma_8^2/\mathrm{d}\ln k}
= 4.8\times10^{-4} \;=\; 0.048\%.
\end{equation}
A function that is identically unity below $k_c$ therefore cannot alter $\sigma_8$,
whatever it does above $k_c$. This is a statement about filters, not about physics,
and it holds for any modification of this form.

\begin{figure}[htb]
\centering
\includegraphics[width=0.78\textwidth]{figures/fig1_window.png}
\caption{\textbf{Why the baseline cannot work.} \emph{Top:} the frozen transfer
function and the corresponding power suppression $T^2$. \emph{Bottom:} the
contribution per logarithmic interval to $\sigma_8^2$ in $\LCDM$. The modification
switches on exactly where the $\sigma_8$ filter has stopped listening. Only
$0.048\%$ of the variance lies above $k_c$.}
\label{fig:window}
\end{figure}

%---------------------------------------------------------------------
\subsection{Correction 2: \texorpdfstring{$S_8$}{S8} and \texorpdfstring{$\sigma_8$}{sigma8} are the same measurement}

Since $\Omega_m$ is fixed by Table~\ref{tab:cosmo} and is unaffected by a transfer
function applied at $z=0$, Eq.~\eqref{eq:S8} is a fixed multiplication:
\begin{equation}
S_8 = 1.02501\,\sigma_8 .
\end{equation}
There is no freedom for one to pass while the other fails. Quantitatively, the
baseline gives
\begin{equation}
\sigma_8=0.8112,\qquad S_8=0.8315,
\end{equation}
against targets $0.76\pm0.03$ and $0.78\pm0.03$, which are themselves consistent
with each other under the same rescaling ($0.76\times1.025=0.779$). The deviations
are
\begin{equation}
\frac{0.8112-0.76}{0.03}=1.71\sigma,
\qquad
\frac{0.8315-0.78}{0.03}=1.72\sigma .
\end{equation}
These are the same number. The earlier annotation of $S_8$ as a pass and $\sigma_8$
as a failure arose from comparing a slightly low quoted value ($0.829$ rather than
the correct $0.8315$) against the looser of two equivalent targets. Both quantities
are $1.7\sigma$ high; neither passes. Figure~\ref{fig:s8} shows the geometry.

\begin{figure}[htb]
\centering
\includegraphics[width=0.66\textwidth]{figures/fig2_s8_sigma8.png}
\caption{\textbf{The claimed pass/fail split is not attainable.} With $\Omega_m$
fixed, the model is confined to the black line. The baseline point coincides with
$\LCDM$ to four decimal places and lies outside both target bands (green). Dotted
guides show how projecting the same point onto the two axes produced the
appearance of one pass and one failure.}
\label{fig:s8}
\end{figure}

%---------------------------------------------------------------------
\subsection{Correction 3: the Lyman-\texorpdfstring{$\alpha$}{alpha} ratio is a convention artefact}

A transfer function multiplies the density contrast, so the power spectrum carries
$T^2$. Evaluating the four-point mean under both conventions:
\begin{equation}
\left\langle \frac{T^2P}{P}\right\rangle_{4} = 0.8125,
\qquad
\left\langle \frac{TP}{P}\right\rangle_{4} = 0.8933 .
\end{equation}
The reported value, $0.891$, matches the second. The quoted Lyman-$\alpha$ pass
therefore rests on applying $T$ to the power spectrum rather than to the density
field, which is off by one power.

A second issue compounds it. Two of the four sampling points, $k=0.5$ and
$k=1.0\,\hMpc$, lie below $k_c$, where $T\equiv1$ by construction. Half the average
is a tautology, and it pulls the mean upward by design. Sampling the range that the
Lyman-$\alpha$ forest actually constrains, under the standard convention, gives
Table~\ref{tab:lya}: a $25\%$ suppression at $k=2$, $50\%$ at $k=3$, $70\%$ at $k=5$,
and $85\%$ at $k=10\,\hMpc$. Suppressions of this size across this range are excluded
by a wide margin; typical analyses tolerate at most tens of percent. See
Figure~\ref{fig:lya}.

\begin{table}[htb]
\centering
\small
\begin{tabular}{lcccccc}
\toprule
$k\;[\hMpc]$ & 1.0 & 2.0 & 3.0 & 5.0 & 8.0 & 10.0\\
\midrule
$P_{\rm mod}/P_{\LCDM}$, standard ($T^2$) & $1.000$ & $0.750$ & $0.500$ & $0.300$ & $0.188$ & $0.150$\\
suppression & $0\%$ & $25\%$ & $50\%$ & $70\%$ & $81\%$ & $85\%$\\
\bottomrule
\end{tabular}
\caption{Power suppression from the frozen $T(k)$ under the standard convention,
over the range the Lyman-$\alpha$ forest constrains. The quoted ratio of $0.891$
describes none of this.}
\label{tab:lya}
\end{table}

\begin{figure}[htb]
\centering
\includegraphics[width=0.72\textwidth]{figures/fig3_lyman_alpha.png}
\caption{\textbf{Two conventions, two conclusions.} The dashed curve is what was
evaluated; the solid curve is the standard $P\propto T^2$. Open circles mark the four
sampling points behind the quoted ratio, two of which sit in the region where
$T\equiv1$ by construction. The green band indicates the order of magnitude of
tolerable suppression.}
\label{fig:lya}
\end{figure}

%---------------------------------------------------------------------
\subsection{Corrected baseline}

\begin{table}[htb]
\centering
\small
\begin{tabular}{lccl}
\toprule
Observable & Previously reported & Verified here & Verdict\\
\midrule
$\sigma_8$ & $0.811$, \textsc{fail} & $0.8112$ & value correct; but it is the
$\LCDM$ value\\
$S_8$ & $0.829$, \textsc{pass} & $0.8315$ & $1.7\sigma$ high; not a pass\\
Lyman-$\alpha$ ratio & $0.891$, \textsc{pass} & $0.8125$ & wrong convention;
excluded under the right one\\
$\Delta\sigma_8/\sigma_8$ & --- & $-4.8\times10^{-5}$ & the mechanism does nothing\\
\bottomrule
\end{tabular}
\caption{Corrected Layer-1 baseline. The frozen $T(k)$ leaves $\sigma_8$ and $S_8$
at their $\LCDM$ values and violates Lyman-$\alpha$ constraints once the standard
convention is used.}
\label{tab:corrected}
\end{table}

\noindent
The honest summary of Layer 1 is that the frozen function has the worst of both
outcomes: it does nothing where help is needed and too much where it is not.

%=====================================================================
\section{The effective coupling from the published Lagrangian}\label{sec:geff}

Layer 1 is already fatal for the observational claim. Layer 2 is worth testing
anyway, because the outcome determines whether any refitted version of the
mechanism could work, and because it governs whether the proposed
six-dimensional programme is worth starting.

\subsection{Action}

The published equation of motion,
\begin{equation}
\ddot r+3H\dot r-\frac{1}{a^2}\nabla^2r+m_r^2r+(\gamma+3\beta_3)r^2+\delta r^3
+\alpha r\!\left(\dot r^2-(\nabla r)^2\right)
=\frac{\lambda}{M_5^{3/2}}(\partial_\mu\phi)^2,
\label{eq:eom}
\end{equation}
follows from
\begin{equation}
S=\int\!\mathrm{d}^4x\sqrt{-g}\left[\frac{\Mpl^2}{2}R-\frac12(\partial r)^2-V(r)
-\frac{\alpha}{2}r(\partial r)^2+\frac{\lambda}{M_5^{3/2}}\,r\,(\partial\phi)^2\right]
+S_m[\psi,g_{\mu\nu}],
\label{eq:action}
\end{equation}
with $V(r)=\tfrac12 m_r^2r^2+\tfrac13(\gamma+3\beta_3)r^3+\tfrac14\delta r^4$. Two
features of Eq.~\eqref{eq:action} govern everything below. The radion kinetic term
has standard sign, so the field is not a ghost. And the matter sector $S_m$ is
minimally coupled: there is no conformal factor multiplying $g_{\mu\nu}$ and no
$r\,T^{\mu}{}_{\mu}$ term. The only route from $r$ to matter is the last term in the
bracket, which is of the electromagnetic $F^2$ type.

\subsection{How an \texorpdfstring{$F^2$}{F2} coupling reaches non-relativistic matter}
\label{sec:damour}

Free radiation has $\langle F_{\mu\nu}F^{\mu\nu}\rangle=0$, so the coupling does not
act on the photon background. It reaches ordinary matter through the electromagnetic
self-energy of bound systems, principally the Coulomb energy of nuclei. This is the
standard construction for dilaton- and moduli-like fields. Writing $\xi_A$ for the
fraction of the mass of body $A$ contributed by electromagnetic binding,
\begin{equation}
m_A(r)=m_A\left[1+\xi_A\frac{\lambda}{M_5^{3/2}}\,r+\mathcal{O}(r^2)\right],
\qquad \xi_A\sim10^{-3},
\end{equation}
and the dimensionless coupling that controls all scalar-mediated effects is
\begin{equation}
\boxed{\;\beta_A \equiv \Mpl\frac{\partial\ln m_A}{\partial r}
= \Mpl\,\xi_A\,\frac{\lambda}{M_5^{3/2}}\;}
\label{eq:beta}
\end{equation}
Two properties matter later. First, $\beta_A$ is independent of the field value,
because the coupling in Eq.~\eqref{eq:action} is linear in $r$. Second, $\xi_A$
depends on composition, roughly as $Z^2/A^{4/3}$, so $\beta_A$ differs between
materials.

\subsection{Quasi-static reduction}

Write $r=\bar r(t)+\delta r(\mathbf{k},t)$ and keep the action to quadratic order.
The cubic and quartic terms of $V$, and the derivative self-interaction
$\alpha r(\partial r)^2$, contribute to the linear equation for $\delta r$ only
through the background-dependent effective mass
\begin{equation}
\meff^2 = V''(\bar r)\big/\!\left(1+\alpha\bar r\right) + \mathcal{O}(\dot{\bar r}^2),
\end{equation}
and do not generate any new coupling between $\delta r$ and $\delta\rho_m$ at this
order. (This is checked explicitly in Appendix~A.) Inside the sound horizon and for
$k/a\gg H$ the time derivatives of $\delta r$ are negligible against $k^2/a^2$, and
\begin{equation}
\left(\frac{k^2}{a^2}+\meff^2\right)\delta r
=\frac{\beta}{\Mpl}\,\delta\rho_m
\qquad\Longrightarrow\qquad
\delta r=\frac{\beta}{\Mpl}\,
\frac{\delta\rho_m}{k^2/a^2+\meff^2}.
\label{eq:deltar}
\end{equation}
Matter responds to the gradient of $\Phi+(\beta/\Mpl)\delta r$, with the Newtonian
potential obeying $-k^2\Phi/a^2=4\pi G\,\delta\rho_m$. Substituting
Eq.~\eqref{eq:deltar} into the geodesic equation and using $8\pi G\Mpl^2=1$,
\begin{equation}
\ddot\delta_m+2H\dot\delta_m=4\pi G_{\rm eff}(k,a)\,\rho_m\delta_m,
\end{equation}
\begin{equation}
\boxed{\;
\frac{\Geff(k,a)}{G}=1+2\beta^2\,\frac{k^2}{k^2+a^2\meff^2}\;}
\label{eq:geff}
\end{equation}
This is the familiar scalar-tensor result, and its familiarity is the point: nothing
about the braneworld origin of $r$ modifies it.

%=====================================================================
\section{Three independent failures}\label{sec:failures}

\subsection{Sign}

The correction in Eq.~\eqref{eq:geff} is proportional to $\beta^2$. It is positive
for every real value of $\lambda$, $M_5$ and $\xi_A$, and for either sign of the
coupling. This is not an artefact of the parameterisation. Exchange of a spin-zero
quantum between two sources of like charge is attractive, so any healthy scalar
linearly coupled to matter increases the effective gravitational attraction. There
is no choice of parameters for which $\Geff<G$.

The consequence is direct: \emph{the mechanism as written raises $\sigma_8$ and
$S_8$}. It moves the observables in the wrong direction. Integrating the growth
equation numerically (Appendix~A) gives, for $\beta=0.10$,
$\Delta P/P=+16\%$ at $k=5\,\hMpc$, and for $\beta=0.22$, $+104\%$.
Figure~\ref{fig:geff} shows both panels.

The only formal escape is a wrong-sign kinetic term for $r$, which would make the
scalar a ghost. That is not a physical option and we do not pursue it.

\begin{figure}[htb]
\centering
\includegraphics[width=0.96\textwidth]{figures/fig4_geff.png}
\caption{\textbf{The derived response, against the response required.}
\emph{(a)} $\Geff/G$ from Eq.~\eqref{eq:geff}: above unity everywhere, rising through
the transition at $k=a\meff$, saturating at $1+2\beta^2$.
\emph{(b)} The resulting fractional change in $P(k)$ at $z=0$ (red) against the
change the fitted $T(k)$ encodes (blue). They occupy opposite half-planes.}
\label{fig:geff}
\end{figure}

\subsection{Shape}

Even setting sign aside, the functional forms are incompatible. Equation
\eqref{eq:geff} describes a \emph{saturating step}: flat and negligible for
$k\ll a\meff$, rising through a transition, and constant at $1+2\beta^2$ for
$k\gg a\meff$. The fitted Eq.~\eqref{eq:Tk} is the mirror image: flat at unity for
$k<k_c$, then declining without bound as $k^{-1/2}$.

The distinction is not cosmetic. The Lyman-$\alpha$ behaviour of the fitted form
depends entirely on continued decline in the tail. The dynamics gives a plateau
there. No value of $\meff$, $\beta$ or $M_5$ interpolates between a bounded step and
an unbounded power law.

A further mismatch appears in panel (b) of Figure~\ref{fig:geff}. The predicted
$\Delta P/P$ is non-zero even at $k=0.01\,\hMpc$, where one might expect the Yukawa
cutoff to protect large scales. It is not protected, because the comoving Compton
scale $a\meff$ shrinks towards zero as $a\to0$: at early times the force is
effectively long-range at all comoving wavenumbers, and every mode accumulates some
enhancement. The result is closer to a scale-independent amplitude shift than to a
localised feature. We return to this in Section~\ref{sec:comoving}.

\subsection{Magnitude}

Reaching $\sigma_8=0.76$ requires $|\Delta P/P|\sim0.1$ near $k\sim0.2\,\hMpc$, hence
$\beta\simeq0.2$ and $\beta^2\simeq5\times10^{-2}$.

The available coupling is far smaller. Because $\meff^{-1}\sim$ Mpc, the force is of
effectively infinite range for every laboratory and Solar-System test, and because
the coupling is linear in $r$ there is no screening (Section~\ref{sec:screening}).
Two independent classes of bound apply.

\begin{itemize}[leftmargin=1.4em,itemsep=2pt]
\item \textbf{Composition-independent.} A long-range scalar coupled to matter with
strength $\beta$ shifts the PPN parameter $\gamma$ by $\Delta\gamma\simeq-2\beta^2$.
The Cassini bound $|\gamma-1|<2.3\times10^{-5}$ then gives
$\beta^2\lesssim10^{-5}$, a shortfall of at least four orders of magnitude. We treat
this as the conservative floor.
\item \textbf{Composition-dependent.} Because $\xi_A$ varies between materials,
$\beta_A$ does too, and the scalar produces a differential acceleration.
Schematically $\eta\sim\beta^2\,\Delta\xi/\bar\xi$ with $\Delta\xi/\bar\xi=\mathcal{O}(0.1$--$1)$
between the Ti and Pt test masses. The MICROSCOPE limit $|\eta|<10^{-15}$ then
implies $\beta^2\lesssim10^{-14}$, a shortfall of thirteen orders.
\end{itemize}

\noindent
The second estimate is schematic and should be replaced by a proper calculation with
tabulated nuclear $\xi_A$ before it is quoted as a number. The first does not depend
on that refinement. Either way the requirement is not close to being met; see
Figure~\ref{fig:coupling}.

\begin{figure}[htb]
\centering
\includegraphics[width=0.84\textwidth]{figures/fig5_coupling.png}
\caption{\textbf{Required versus permitted coupling.} The red band is what a
percent-level effect on $\sigma_8$ demands. The blue band is the conservative
Solar-System floor; the green band is the nominal equivalence-principle limit. The
gap is at least four orders of magnitude and plausibly much larger.}
\label{fig:coupling}
\end{figure}

\begin{keybox}
\textbf{Verdict on Layer 2.} The published Lagrangian produces an enhancement rather
than a suppression, with a saturating rather than declining $k$-dependence, at a
coupling strength that existing experiments already exclude by at least four orders
of magnitude. The attribution of the fitted $T(k)$ to radion leakage is falsified,
not merely underived.
\end{keybox}

%=====================================================================
\section{Withdrawal of the screening result}\label{sec:screening}

An earlier calculation in this project reported a ratio of $10^{42}$ between local
and cosmological values of $(\partial_\mu\phi)^2$, and interpreted it as a screening
failure of that magnitude. We withdraw that interpretation.

The quantity computed is the ratio of \emph{background source strengths}, which sets
the local value of $\bar r$ through Eq.~\eqref{eq:eom}. It is not a screening factor.
Screening means that the \emph{coupling} $\beta$ is suppressed in dense
environments, which happens in three known ways: a density-dependent mass
(chameleon), a density-dependent coupling (symmetron), or derivative
self-interaction that suppresses the scalar gradient near a source (Vainshtein).
Equation~\eqref{eq:action} realises none of them. The coupling is linear in $r$,
so $\beta$ in Eq.~\eqref{eq:beta} is a constant, and the $\alpha r(\partial r)^2$
term acts on the radion's own gradient, not on its coupling to matter.

The corrected statement is stronger than the one withdrawn, and in the same
direction:

\begin{keybox}
There is no screening mechanism in this Lagrangian. The composition dependence of
$\xi_A$ therefore propagates without suppression into an Eötvös signal, and the
magnitude bound of Section~\ref{sec:failures} applies with full force.
\end{keybox}

%=====================================================================
\section{The comoving-scale problem}\label{sec:comoving}

A feature at fixed \emph{comoving} wavenumber cannot be produced by a fixed proper
length. Compactification radii, Compton wavelengths and force ranges are proper
lengths; $k_c$ in Eq.~\eqref{eq:Tk} is comoving. Every genuine feature in the
observed power spectrum arises instead from a time-dependent crossing condition:
the horizon size at matter--radiation equality, a free-streaming length, or a Jeans
scale.

For a scalar mediator the natural identification is $k_c\simeq a\,\meff$, evaluated
today. With $k_c=1.5\,\hMpc=1.01\,\mathrm{Mpc}^{-1}$ and
$1\,\mathrm{Mpc}^{-1}=6.39\times10^{-30}$ eV,
\begin{equation}
\meff \simeq 6.5\times10^{-30}\ \mathrm{eV}
\;\simeq\; 4.5\times10^{3}\,H_0,
\qquad
\meff^{-1}\simeq1\ \mathrm{Mpc}.
\label{eq:mass}
\end{equation}
This is the number any successor construction must produce. It is worth recording
how far it sits from a generic braneworld modulus mass. If $\meff$ is of order the
Kaluza--Klein scale, sub-millimetre tests of Newtonian gravity require
$\meff\gtrsim10^{-2}$~eV, twenty-eight orders of magnitude too heavy. If instead the
modulus is a classically flat direction lifted only by whatever generates the
residual vacuum energy, one expects $\meff\sim H_0$, and Eq.~\eqref{eq:mass} then
demands a factor of a few thousand above that. The second regime is the only one
within reach, and it is not generic.

The same argument has a second edge. If $\meff$ is constant, the comoving feature
drifts as $a^{-1}$, so the correct object is $T(k,z)$, not $T(k)$. The frozen
baseline has no redshift dependence at all. Fitting a $z$-independent function to
data drawn from several epochs is therefore inconsistent with the mechanism it was
meant to represent, independently of everything in
Sections~\ref{sec:geff}--\ref{sec:failures}.

%=====================================================================
\section{What a viable successor would require}\label{sec:requirements}

Collecting the constraints, any mechanism that still wishes to be described as
radion leakage and to suppress rather than enhance small-scale power must satisfy
all of the following.

\begin{enumerate}[leftmargin=1.6em,itemsep=4pt]
\item \textbf{$\Delta G<0$, which rules out conformal scalar exchange.} The two
known routes are genuine energy transfer out of the matter sector, so that
$\dot\rho_m+3H\rho_m=-\Gamma\rho_m$ with $\Gamma$ derived rather than inserted, or a
derivative or disformal coupling in the Horndeski class, where $\Geff<G$ is
attainable in some parameter ranges. The first is what the word ``leakage'' ought to
mean and is \emph{not} what Eq.~\eqref{eq:action} describes. Either route requires
rewriting the coupling structure, not retuning it.

\item \textbf{A suppression that does not drift.} Either $\meff\propto H$, in which
case the feature tracks the horizon and will not sit at $1.5\,\hMpc$, or the model
predicts an explicit $T(k,z)$ that must be fitted jointly across redshift.

\item \textbf{A coupling that avoids the electromagnetic sector, or is screened.}
Given Section~\ref{sec:failures}, an unscreened $F^2$ coupling of the required
strength is excluded. A chameleon or Vainshtein sector would have to be added
explicitly and shown to be stable.

\item \textbf{Support where $\sigma_8$ lives.} Any $T$ that equals unity below
$1\,\hMpc$ cannot move $\sigma_8$, by Figure~\ref{fig:window}. This is a constraint
on the fitting function itself and is independent of the mechanism.

\item \textbf{Simultaneous consistency with Lyman-$\alpha$.} Conditions 4 and 5
together are the original difficulty, and nothing in this analysis has made it
easier.
\end{enumerate}

%=====================================================================
\section{Implications for the six-dimensional programme}\label{sec:6d}

A move to a six-dimensional codimension-two construction, stabilised by bulk
magnetic flux and a bulk cosmological constant, has been proposed as a way to derive
$k_c$ from geometry rather than fitting it. We recommend deferring it, for three
reasons.

First, none of the failures in Section~\ref{sec:failures} depends on the number of
dimensions. Equation~\eqref{eq:geff} follows from a light scalar with a healthy
kinetic term coupled linearly to matter through $F^2$. That description survives
dimensional reduction from five, six or ten dimensions. Flux compactification would
determine how $\meff$ and $\lambda$ are fixed; it would not change the sign of
$\beta^2$.

Second, the claimed advantages of codimension-two are weaker than advertised. In the
football and rugby-ball solutions the bulk flux must be tuned to obtain flat branes,
the compact case requires two branes of exactly equal tension, and the apparent
insensitivity of the brane geometry to brane vacuum energy does not survive a change
of tension: these models remain subject to Weinberg's no-go theorem. Flux
quantisation discretises the compactification radius; it does not select it, since
the radius is fixed by a balance among the bulk cosmological constant, the gauge
coupling and the flux integer. One fitted number would be traded for three inputs.

Third, the premise that five dimensions cannot achieve dynamical scale selection is
not correct. Goldberger--Wise stabilisation generates an exponential hierarchy from
$\mathcal{O}(1)$ bulk parameters. The obstacle identified in
Section~\ref{sec:comoving} is not the number of dimensions but the mapping from any
proper length to a comoving wavenumber, and adding a dimension does not address it.

The refined claim that does survive is narrower: codimension-two is the lowest
setting in which the modulus potential comes from bulk field equations plus a
topological quantisation condition, rather than from an added stabilising scalar.
That is an improvement in economy. It is not an argument that the resulting scale is
$1.5\,\hMpc$.

%=====================================================================
\section{Falsifiable content}\label{sec:predictions}

This paper makes three statements that can be checked directly and would overturn
the conclusion if found wrong.

\begin{enumerate}[leftmargin=1.6em,itemsep=3pt]
\item \textbf{The window-function claim.} Any Boltzmann code, with any reasonable
cosmology, will find that a transfer function identically unity below $1.5\,\hMpc$
changes $\sigma_8$ by less than $10^{-4}$ fractionally. This requires no model
assumptions and can be confirmed in minutes with \texttt{verify\_baseline.py}.
\item \textbf{The sign of $\Geff$.} If an explicit dimensional reduction of the 5D
Gauss--Bonnet action yields a matter coupling that is \emph{not} of the conformal
form assumed in Section~\ref{sec:damour}, for example a genuine disformal or
derivative coupling, then Eq.~\eqref{eq:geff} does not apply and the sign
argument must be redone. Producing that reduction explicitly is the natural
next calculation, and is the one route by which the mechanism could be revived.
\item \textbf{The magnitude bound.} A careful Damour--Donoghue treatment with
tabulated $\xi_A$ would sharpen the equivalence-principle constraint. We predict it
will tighten, not relax, the bound quoted in Section~\ref{sec:failures}.
\end{enumerate}

\noindent
The prediction that would have been the interesting one, a distinctive
$\Delta P(k,z)$ fingerprint testable against forthcoming wide-field data, is not
available: the mechanism predicts enhancement at a level already excluded, so there
is nothing left to test against future surveys.

%=====================================================================
\section{Conclusions}

The broken power-law $T(k)=(1.5/k)^{1/2}$ remains a well-defined object that can be
compared with data. What it does when compared is not what was reported. It leaves
$\sigma_8$ and $S_8$ at their $\LCDM$ values, because $99.95\%$ of the $\sigma_8$
variance lies below its break scale, and it suppresses the matter power spectrum by
$50$--$85\%$ across the Lyman-$\alpha$ range once the standard $P\propto T^2$
convention is applied. The reported pattern of an $S_8$ pass alongside a $\sigma_8$
failure is not attainable at all, since the two quantities differ by a fixed factor
of $1.025$.

The attribution of the fitted curve to a radion in a Gauss--Bonnet braneworld is
separately falsified. The published Lagrangian yields
$\Geff/G=1+2\beta^2k^2/(k^2+a^2\meff^2)$: an enhancement, of saturating rather than
power-law shape, at a coupling that Solar-System and equivalence-principle tests
already exclude by at least four orders of magnitude. The earlier screening
calculation is withdrawn, and replaced by the stronger statement that the Lagrangian
contains no screening mechanism at all.

We do not regard this as a reason to abandon the underlying programme, but we do
regard it as a reason to change what is being built. The five-dimensional background
numerics and the reproducibility infrastructure developed for this project are
sound and unaffected. The next useful calculation is an explicit dimensional
reduction of the 5D action, carried far enough to determine whether the induced
matter coupling is conformal, as assumed here, or of a derivative type that could
evade the sign argument. Until that is done, neither a refit of the transfer
function nor a move to six dimensions rests on anything.

%=====================================================================
\section*{Data and code availability}
\addcontentsline{toc}{section}{Data and code availability}

All scripts, outputs and figures accompanying this paper are released with it:
\texttt{verify\_baseline.py} (independent \textsc{camb} run and observable
extraction), \texttt{make\_figures.py} (all figures),
\texttt{baseline\_verification.json} (numerical results),
\texttt{pk\_baseline.csv} (power spectra and the $\sigma_8$ integrand), and
\texttt{V2.0\_BASELINE\_LOCKED.md} (the frozen definitions of
Section~\ref{sec:baseline}). Every number in this paper is regenerated by running
the two scripts in order. Repository:
\url{https://github.com/GeometricCosmo/gb-leakage-cmb}.

\section*{Acknowledgements}
\addcontentsline{toc}{section}{Acknowledgements}

The analysis in Sections~\ref{sec:verify}--\ref{sec:failures} was carried out with
substantial assistance from large language model tooling, used for derivation
checking, code generation and adversarial review. All numerical results were
executed and are reproducible from the released scripts. Responsibility for the
content, including any remaining errors, rests with the author.

\newpage
\appendix

%=====================================================================
\section{Growth equation and the role of the higher-order terms}

\subsection{Numerical growth integration}

Writing $N=\ln a$ and using $\mathrm{d}\ln H/\mathrm{d}N=-\tfrac32\Omega_m(a)$ for a
flat $\Lambda$ background,
\begin{equation}
\frac{\mathrm{d}^2\delta}{\mathrm{d}N^2}
+\left(2-\tfrac32\Omega_m(a)\right)\frac{\mathrm{d}\delta}{\mathrm{d}N}
-\tfrac32\Omega_m(a)\,\frac{\Geff(k,a)}{G}\,\delta = 0,
\qquad
\Omega_m(a)=\frac{\Omega_m a^{-3}}{\Omega_m a^{-3}+1-\Omega_m}.
\end{equation}
We integrate from $a=10^{-3}$ with $\delta=\mathrm{d}\delta/\mathrm{d}N=a$ and form
$\Delta P/P=(\delta/\delta_{\LCDM})^2-1$ at $a=1$. Results are in
Table~\ref{tab:growth}.

\begin{table}[htb]
\centering
\small
\begin{tabular}{lcccccc}
\toprule
$k\,[\hMpc]$ & $0.01$ & $0.10$ & $0.50$ & $1.50$ & $5.0$ & $20$\\
\midrule
$\Delta P/P$, $\beta=0.10$ & $+0.047$ & $+0.105$ & $+0.143$ & $+0.158$ & $+0.161$ & $+0.162$\\
$\Delta P/P$, $\beta=0.22$ & $+0.244$ & $+0.612$ & $+0.891$ & $+1.007$ & $+1.039$ & $+1.043$\\
\bottomrule
\end{tabular}
\caption{Fractional change in $P(k)$ at $z=0$ from Eq.~\eqref{eq:geff}, with
$\meff=1.5\,\hMpc$. Positive at every scale, monotonically increasing, saturating.}
\label{tab:growth}
\end{table}

\subsection{Why the cubic, quartic and derivative terms drop out}

Expanding Eq.~\eqref{eq:eom} to first order in $\delta r$ about a homogeneous
background $\bar r(t)$:
\begin{align}
\text{from } (\gamma+3\beta_3)r^2 :&\quad 2(\gamma+3\beta_3)\bar r\,\delta r,\\
\text{from } \delta r^3 :&\quad 3\delta\,\bar r^2\,\delta r,\\
\text{from } \alpha r(\dot r^2-(\nabla r)^2) :&\quad
\alpha\dot{\bar r}^2\,\delta r + 2\alpha\bar r\dot{\bar r}\,\dot{\delta r}.
\end{align}
The first two are proportional to $\delta r$ with time-dependent coefficients and are
absorbed into $\meff^2=V''(\bar r)$. The third contributes a friction-like term
$2\alpha\bar r\dot{\bar r}\,\dot{\delta r}$, which is negligible in the quasi-static
limit where $\dot{\delta r}$ is dropped relative to $(k^2/a^2)\delta r$, and a mass
shift $\alpha\dot{\bar r}^2$. None of them couples $\delta r$ to $\delta\rho_m$.
The matter coupling therefore comes entirely from the $F^2$ term, as used in
Eq.~\eqref{eq:deltar}.

The one caveat worth recording: if $\bar r$ evolves rapidly, so that
$\alpha\bar r\dot{\bar r}$ is not small compared with $H$, the quasi-static
approximation itself requires re-examination. For the parameter values considered
here $\bar r$ is slowly rolling and the approximation holds, but a construction that
relied on fast radion evolution would need this checked.

%=====================================================================
\section{Reproduction}

\begin{verbatim}
pip install camb numpy scipy matplotlib
python verify_baseline.py     # -> baseline_verification.json, pk_baseline.csv
python make_figures.py        # -> figures/fig1..fig5 .png
\end{verbatim}

\noindent
Runtime is under two minutes on a single core. \textsc{camb} version 2.0.4 was used;
any version from 1.3 onward should reproduce $\sigma_8$ to better than $10^{-3}$.
The quadrature in Eq.~\eqref{eq:sigma8} is Simpson's rule on $4000$ logarithmic
points over $10^{-4}$--$100\,\hMpc$; halving the resolution changes $\sigma_8$ in the
fifth decimal.

%=====================================================================
\section{Relationship to earlier versions}

\begin{table}[htb]
\centering
\small
\begin{tabular}{lll}
\toprule
Version & Principal claim & Status after this paper\\
\midrule
v1.8.2 & Exponential $T(k)$ resolves small-scale abundances & superseded by v1.9.x\\
v1.8.3 & $\sigma_8=0.76$ from the stated $T(k)$ & incorrect; not reproducible\\
v1.9.0 & $S_8$/Lyman-$\alpha$ mutually exclusive; mechanism excluded & directionally correct\\
v1.9.1 & Power-law $T(k)$ passes $S_8$ and Lyman-$\alpha$ & \textbf{retracted} (Sec.~\ref{sec:verify})\\
v2.0 & Baseline fails; mechanism falsified & this paper\\
\bottomrule
\end{tabular}
\caption{Version history and current status.}
\end{table}

\noindent
The v1.9.0 conclusion, that the mechanism is observationally excluded, was closer to
correct than the v1.9.1 revision that replaced it. The intervening change of sign in
the project's own conclusions traces to the two convention errors identified in
Section~\ref{sec:verify}.

%=====================================================================
\section{Errata}\label{app:errata}

The following specific statements in earlier releases should not be cited.

\begin{enumerate}[leftmargin=1.6em,itemsep=3pt]
\item ``$S_8=0.829$ passes the weak-lensing constraint.'' It does not; the correct
value is $0.8315$ and it lies $1.7\sigma$ above the target.
\item ``The mechanism passes $S_8$ while failing $\sigma_8$.'' Not attainable;
$S_8=1.025\,\sigma_8$ identically.
\item ``Lyman-$\alpha$ ratio $=0.891$, within bounds.'' Computed with
$P_{\rm mod}=TP$ rather than $T^2P$, and averaged over sampling points at which
$T\equiv1$. The correct four-point value is $0.8125$ and the suppression reaches
$70\%$ at $k=5\,\hMpc$.
\item ``Gradient screening fails by $10^{42}$.'' Category error; withdrawn. See
Section~\ref{sec:screening}.
\item ``The holographic derivation fixes $T(k)$ uniquely.'' The asymptotic index has
holographic motivation; the break scale $k_c$ and the low-$k$ behaviour are fitted.
\end{enumerate}

%=====================================================================
\begin{thebibliography}{99}
\small
\bibitem{planck18} Planck Collaboration, \emph{Planck 2018 results. VI.
Cosmological parameters}, Astron.\ Astrophys.\ \textbf{641} (2020) A6
[arXiv:1807.06209].
\bibitem{camb} A.~Lewis, A.~Challinor and A.~Lasenby, \emph{Efficient computation of
CMB anisotropies in closed FRW models}, Astrophys.\ J.\ \textbf{538} (2000) 473
[astro-ph/9911177].
\bibitem{damour94} T.~Damour and A.~M.~Polyakov, \emph{The string dilaton and a
least coupling principle}, Nucl.\ Phys.\ B \textbf{423} (1994) 532
[hep-th/9401069].
\bibitem{damour10} T.~Damour and J.~F.~Donoghue, \emph{Equivalence principle
violations and couplings of a light dilaton}, Phys.\ Rev.\ D \textbf{82} (2010)
084033 [arXiv:1007.2792].
\bibitem{microscope} P.~Touboul et al.\ (MICROSCOPE Collaboration), \emph{MICROSCOPE
mission: final results of the test of the equivalence principle}, Phys.\ Rev.\ Lett.\
\textbf{129} (2022) 121102 [arXiv:2209.15487].
\bibitem{cassini} B.~Bertotti, L.~Iess and P.~Tortora, \emph{A test of general
relativity using radio links with the Cassini spacecraft}, Nature \textbf{425}
(2003) 374.
\bibitem{carroll03} S.~M.~Carroll and M.~M.~Guica, \emph{Sidestepping the
cosmological constant with football-shaped extra dimensions}, hep-th/0302067.
\bibitem{navarro03} I.~Navarro, \emph{Codimension two compactifications and the
cosmological constant problem}, JCAP \textbf{0309} (2003) 004 [hep-th/0302129].
\bibitem{garriga04} J.~Garriga and M.~Porrati, \emph{Football shaped extra
dimensions and the absence of self-tuning}, JHEP \textbf{0408} (2004) 028
[hep-th/0406158].
\bibitem{akerblom10} N.~Akerblom and G.~Cornelissen, \emph{A compact codimension-two
braneworld with precisely one brane}, Phys.\ Rev.\ D \textbf{81} (2010) 124025
[arXiv:1004.1807].
\bibitem{gw99} W.~D.~Goldberger and M.~B.~Wise, \emph{Modulus stabilization with
bulk fields}, Phys.\ Rev.\ Lett.\ \textbf{83} (1999) 4922 [hep-ph/9907447].
\bibitem{horndeski} G.~W.~Horndeski, \emph{Second-order scalar-tensor field
equations in a four-dimensional space}, Int.\ J.\ Theor.\ Phys.\ \textbf{10} (1974)
363.
\bibitem{adelberger09} E.~G.~Adelberger, J.~H.~Gundlach, B.~R.~Heckel, S.~Hoedl and
S.~Schlamminger, \emph{Torsion balance experiments: a low-energy frontier of
particle physics}, Prog.\ Part.\ Nucl.\ Phys.\ \textbf{62} (2009) 102.
\bibitem{zenodo191} A.~Swart, \emph{Radion leakage in a 5D Gauss--Bonnet braneworld},
Zenodo v1.9.1 (2026), \url{https://zenodo.org/records/21663023}.
\end{thebibliography}

\end{document}
